% Clase 16 --- Por qué la normal se elige con la mano derecha (§1.4).
% "Pareciera que la elección que hicimos de N no es la más conveniente."
% Los dos primeros paneles tienen EXACTAMENTE la misma física (misma espira,
% misma corriente, mismas fuerzas, mismo par): lo único que cambia es cuál de
% las dos normales se elige. Con la elección "a ojo" el equilibrio estable cae
% en theta = 180 y el signo de la fórmula no acompaña; con la de la mano derecha
% cae en theta = 0 y tau = I A x B da módulo, dirección y sentido de una.
\begin{center}
\begin{tikzpicture}[scale=1]

  % =============== (a) normal elegida a ojo ===============
  \begin{scope}
    \draw[figvecaux, figgray] (-2.0,2.35) -- (1.9,2.35);
    \node[figlblaux, anchor=south east] at (1.9,2.4) {$\vec B$};
    \draw[figaux] (0,0) -- (1.5,0);

    \draw[figline, line width=1.6pt] (-1.15,0.80) -- (1.15,-0.80);
    \node[figblue, scale=0.85] at (1.15,-0.80) {$\otimes$};
    \node[figblue, scale=0.85] at (-1.15,0.80) {$\odot$};
    \node[figgray, circle, draw=figgray, fill=white, inner sep=1.4pt] at (0,0) {};

    \draw[figvecres] (1.15,-0.97) -- (1.15,-1.75);
    \node[figlbl, figred, anchor=west] at (1.22,-1.4) {$\vec F_3$};
    \draw[figvecres] (-1.15,0.97) -- (-1.15,1.75);
    \node[figlbl, figred, anchor=east] at (-1.22,1.4) {$\vec F_4$};

    % la normal "a ojo"
    \draw[figvec, line width=1.1pt] (0,0) -- (-0.86,-1.23);
    \node[figlbl, figblue, anchor=north east] at (-0.80,-1.16) {$\hat n$};
    \draw[figaux] (0.95,0) arc (0:-125:0.95);
    \node[figlbl, anchor=west] at (0.58,-1.02) {$\theta$};

    % sentido del par
    \draw[figvecaux, figred, line width=1pt,
          -{Latex[length=2.2mm,width=1.6mm]}] (-1.0,-1.73) arc (240:186:2.0);
    \node[figlblaux, figred, anchor=east] at (-1.95,-1.15) {sentido del par};

    \node[figlbl, anchor=north, align=center] at (0,-2.75)
      {(a) con $\hat n$: el par lleva a\\[-0.2em]
       $\hat n$ \textbf{antiparalelo} a $\vec B$};
  \end{scope}

  % =============== (b) normal por la mano derecha ===============
  \begin{scope}[xshift=4.8cm]
    \draw[figvecaux, figgray] (-2.0,2.35) -- (1.9,2.35);
    \node[figlblaux, anchor=south east] at (1.9,2.4) {$\vec B$};
    \draw[figaux] (0,0) -- (1.5,0);

    \draw[figline, line width=1.6pt] (-1.15,0.80) -- (1.15,-0.80);
    \node[figblue, scale=0.85] at (1.15,-0.80) {$\otimes$};
    \node[figblue, scale=0.85] at (-1.15,0.80) {$\odot$};
    \node[figgray, circle, draw=figgray, fill=white, inner sep=1.4pt] at (0,0) {};

    \draw[figvecres] (1.15,-0.97) -- (1.15,-1.75);
    \node[figlbl, figred, anchor=west] at (1.22,-1.4) {$\vec F_3$};
    \draw[figvecres] (-1.15,0.97) -- (-1.15,1.75);
    \node[figlbl, figred, anchor=east] at (-1.22,1.4) {$\vec F_4$};

    % la normal correcta
    \draw[figvec, line width=1.1pt] (0,0) -- (0.86,1.23);
    \node[figlbl, figblue, anchor=south west] at (0.80,1.16) {$\hat n'$};
    \draw[figaux] (0.95,0) arc (0:55:0.95);
    \node[figlbl, anchor=west] at (0.98,0.42) {$\theta$};

    \draw[figvecaux, figred, line width=1pt,
          -{Latex[length=2.2mm,width=1.6mm]}] (-1.0,-1.73) arc (240:186:2.0);

    \node[figlbl, anchor=north, align=center] at (0,-2.75)
      {(b) con $\hat n'$: el par lleva a\\[-0.2em]
       $\hat n'$ \textbf{paralelo} a $\vec B$};
  \end{scope}

  % =============== (c) la regla, de frente ===============
  \begin{scope}[xshift=9.2cm]
    \draw[figline, line width=1.4pt] (-1.0,-0.85) rectangle (1.0,0.85);
    \draw[figvecres] (0.35,-0.85) -- (1.0,-0.85) -- (1.0,-0.15);
    \draw[figvecres] (-0.35,0.85) -- (-1.0,0.85) -- (-1.0,0.15);
    \node[figlbl, figred, anchor=north] at (0.7,-0.92) {$I$};

    \node[figblue, scale=1.1] at (0,0) {$\odot$};
    \node[figlbl, figblue, anchor=west] at (0.18,0.02) {$\hat n'$};

    \node[figlbl, anchor=north, align=center] at (0,-2.75)
      {(c) dedos $=$ corriente,\\[-0.2em] pulgar $=\hat n'$};
  \end{scope}

\end{tikzpicture}\\[0.5em]
{\small\itshape La física de los paneles (a) y (b) es idéntica ---las fuerzas y
el par son los mismos---; lo único que cambia es cuál de las dos normales se
llama ``la'' normal. Con $\hat n$ el sistema queda con la incomodidad de que la
posición estable es $\theta=180^\circ$, y al escribir el par como producto
vectorial el sentido sale al revés. Eligiendo $\hat n'=-\hat n$ mediante la
regla de la mano derecha ---dedos siguiendo la corriente, pulgar dando la
normal, panel (c)--- la posición estable pasa a ser $\theta=0$ y basta definir
$\vec A = A\,\hat n'$ para que
$\boldsymbol{\tau} = I\,\vec A\times\vec B$
dé de una sola vez módulo, dirección y sentido. El orden importa: $\vec
B\times\vec A$ daría el sentido opuesto, y por eso se escribe $\vec A$ primero.
Con $\boldsymbol{\mu}=I\vec A$ queda $\boldsymbol{\tau}=\boldsymbol{\mu}\times
\vec B$, calcado del dipolo eléctrico $\boldsymbol{\tau}=\vec p\times\vec E$.}
\end{center}
